\documentclass[12pt]{article}
%\usepackage[OT1]{fontenc}
%\usepackage{babel}
%\usepackage{hyperref} 
%\usepackage{url}   %this allows us to cite URLs in the text
\usepackage{graphicx}  %allows for graphic to float when doing jou or doc style
\usepackage{amssymb}  %use formatting tools  for math symbols
\usepackage{natbib}  %might conflict with apa...
\usepackage{booktabs}
%\usepackage{Sweave}
%\usepackage{tikz}
%\usepackage{pgf}
\usepackage{fullpage}
\usepackage{setspace} 
%\usepackage{lineno}
%\linenumbers


\begin{document}  
\title{Research Statement}
\author{Solomon Messing}
\date{October 29, 2014}
\maketitle 

%\begin{center} 
%\Large Research Statement (Messing)
%\end{center}


%\doublespacing 

%I'd start the first paragraph with something about the political science thing.  So, the big thing that you're 
%
%I study how our day-to-day interactions and the strategic communication of political elites affects the democratic process---what citizens learn about politics, how they evaluate political actions, and when they decide to participate.  WHAT IS THE TOP LINE FINDING FROM YOUR WORK?  To understand how social processes shape politics, I use massive data sets from online community, social networks, and political text.  I analyze the data sets using new machine learning techniques, that provide insights.
%
%The rest of the statement needs to emphasize your substantive findings.  Be specific about what you've found and how you find it. Much of what you're saying in this research statement is a summary of what you do, rather than what you find and why it matters and how what you do is ground breaking.  

I study how technological developments in communication affect policymaking and the democratic process---what citizens learn about politics, how they decide to participate, and the actions politicians take and policies they adopt in response.  My work shows that rather than the ``minimal effects'' associated with traditional news media, elite ideas delivered personally shape how people think and how they evaluate policy makers.  I combine experimental interventions, large behavioral data sets on social networks and political text, and statistical and machine learning techniques to generate new insight about how social processes affect the way we think about politics and the policies we adopt.  %I study these processes with a combination of innovative experiments, massive data sets from online communities and political text, analyzed using new techniques from statistics and machine learning to provide insights into public opinion and political behavior.

%New communication and data processing technologies have put personal interactions among individuals---and therefore social influence---at the center of modern political campaigns and constituent communication in America.  These technologies provide tools used by politically active citizens to share discourse with their social contacts, legislators to build support and garner instant feedback from their constituents in a public forum, and political campaigns to message supporters, recruit new activists, raise funds, and build rolodexes.  My research uses the behavioral data generated in the process to study how these micro-interactions coalesce to impact public opinion, political participation, and governance.  %I use experimental interventions; large, multi-source, fine-grained behavioral (``big'') data sets; and statistical and machine learning techniques to examine how social processes affect the way we think about politics.  

%\footnote{I.e., web and mobile data collection, social media APIs, and distributed data systems (especially open-source systems like HBase and Hadoop).}

%My research shows how social processes affect the way we think about politics.  I demonstrate that sharing or endorsing elite opinion changes the way a person's friends respond to political discourse, triggering complex evaluations of the sharer's personal attributes, group membership, and the relationship with the sharer---making it not only more likely that an individual carefully evaluates the discourse, but shaping the \emph{way} she thinks about it.  And, when she encounters more shared political discourse in social media, she is more likely to hold informed opinions, state an interest in politics, and actually vote.  My work also shows that personal communication between representatives and their constituents shapes how individuals evaluate politicians, but does so in ways that politicians exploit to convey influence and cultivate a personal vote.  This work is part of a broader research agenda focused on applying data science to answer substantive questions in political science including substantive questions like how to evaluate the ideology of elites and individuals on the same scale and how to analyze images in political campaigns. This approach has resulted in substantive contributions and external funding in the fields of political science and communication.  

%These technologies affect democratic discourse, levels of political knowledge and the process of representation.  My work shows that structural features of these communication media---their availability, inherent costs, indicators of utility, how closely they link us to our social world, and physical presentation dramatically change the way we consume media.  I marshal large observational datasets, administer ecologically valid experiments and utilize unsupervised machine learning to advance theory about how these new technologies shape outcomes.  

%This approach has resulted in grant awards and substantive contributions in the fields of communication and political science, and I am actively pursuing grants and furthering this research agenda.  

\section*{Dissertation}
My dissertation examines how social processes taking place in a mediated online environment shape the interplay between elite opinion and public opinion, using web experiments and machine learning.  Sharing or endorsing elite opinion changes the way a person's friends respond to political discourse, while doing so in the context of social media's affordances triggering complex evaluations of the sharer's personal attributes, group membership, and the relationship with the sharer.  I show how spreading political discourse this way affects the mix of ideas that an individual encounters and makes it more likely that she carefully evaluates those ideas.  But it also changes the \emph{way} she thinks about the underlying content.  And, when she encounters more shared political discourse in social media, she is more likely to hold informed opinions, state an interest in politics, and actually vote.  


%By examining network structure, media discourse and political attitudes, my dissertation serves to links the literature on information diffusion to individual behavior by accounting for the various social processes that impact both the mix of content to which an individual is exposed and how it is processed.  

%The massive quantity of fine-grained data available thanks to modern web technologies means that for the first time in history, we can study these processes ``in situ.''  My dissertation investigates these phenomena with a series of original web- and field experiments, large-scale behavioral studies of social networks, and analyses of large scale text corpora.  

%After establishing a theoretical foundation for why people share and endorse elite opinion publicly, 

My dissertation first shows that people are more likely to share content from the opinion pages than any other section, meaning that shared content contains not merely facts but persuasive arguments, often about politics (Ch. 1).  I then present an experiment showing that people utilize social media's technological affordances including sharing and endorsement of elite opinion as heuristic cues that serve as easy-to-process signals enabling people to find desirable and/or socially relevant articles with little cognitive burden.  People provided with genuine indications of sharing content (scraped from the web), but not fake indicators, select more articles and evaluate the news reading experience as more enjoyable (Ch. 2).  Next, I show in two experimental studies that social endorsements affect selection to a much greater degree than partisan selective exposure based on source ideology (e.g., Fox News versus MSNBC), and that when social endorsements and partisan source cues are presented simultaneously, partisan selectivity disappears (Ch. 3, co-authored with Sean Westwood and accepted at \emph{Communication Research}).  

My dissertation shows that not all cues are processed equally---when close friends (operationalized based on frequent interaction on Facebook) are randomly assigned to recommend an article, individuals are more likely to read it than when it comes recommended from a peripheral social contact (Ch. 4).  I then show that the relative socioeconomic status and racial group membership of the sharer and the reader drives not only the decision to consume a media item, but also \emph{affects how the reader processes it} (Ch. 5).  

After demonstrating that sharing and endorsing content is a powerful driver of attention, I show that at the aggregate level, this process results in substantial exposure and consumption of discourse that cuts across the ideological divide, by examining every case of sharing elite political opinion during a six-month period in the largest social media platform, Facebook (Ch. 6, co-authored with Eytan Bakshy and Lada Adamic).  Finally, I examine the political implications of these processes by analyzing a 72-day product test with a large sample ($N=1.6M$), which shows that additional exposure to shared political content in social media triggers pro-social political norms, boosts the salience of politics, leads to stronger policy preferences, and increases voter turnout (Ch. 7).  

%My work provides compelling evidence that the social processing of media content triggers complex evaluations of the sharer's relationship with the reader and her social group membership, which affect selectivity and how she processes the content in question.  At the same time, the observational and field experiments triangulate the critical role of social processing in how we think about politics, showing that increased exposure to socially shared and endorsed political content increases the salience of politics, strengthens political opinions, and increases voter turnout.

%We should expect people to engage with political information and opinion differently when an individual or group signals its importance by sharing it or endorsing it.  First and most obviously, sharing a story increases the likelihood it reaches the consumer, which has implications for the mix of content that members of different social strata encounter.  %every model of opinion formation and change is conditional on exposure.   
%Second, importance-signaling affects whether she actually pays attention to it---a hypothesis implied by a wide range of literatures, from surveillance theories suggesting we are hard-wired to monitor our immediate surroundings and especially agents who are socially close to us \cite{lasswell_structure_1960,shoemaker_hardwired_1996}, to dual-processing models of decision making \citep{chaiken_heuristic_1987}, suggesting that we incorporate our peer's judgments when deciding what kind of news to select and attend to \citep{messing_selective_2012}.  This privileged status should increase the likelihood that the consumer will engage deeply with such content and should be expected to make message acceptance/persuasion more likely \citep[e.g.,][]{sherif_attitude_1965,axsom_audience_1987}.  % do we really want to speculate about this, we don't actually measure it.  
%The massive quantity of fine-grained data available thanks to modern web technologies means that for the first time in history, we can study these processes ``in situ.''  %nice

%This research provides a theoretical account of the social processes that occur when individuals endorse elite opinion in a public or social setting. 

\section*{Representation \& Constituent Communication}
These structural changes favoring individual-level interactions have implications beyond public opinion---there are also implications for representative-constituent communication and coalition-building that affect the functioning and responsiveness of government.  One particularly important facet of this communication between representatives and constituents, and what appears to matter most in that realm is claiming credit for bringing particularistic benefits back to a representative's constituents. 

In ``How Words and Money Cultivate a Personal Vote: The Effect of Legislator Credit Claiming on Constituent Credit Allocation'' (with Justin Grimmer and Sean Westwood (2012), \emph{American Political Science Review}), we employ experiments conducted via email and Facebook that dissociated the effect of \emph{claiming credit} for particularistic spending from that spending itself.  Our findings establish that in contrast to theoretical models of credit allocation, increasing the number of statements cultivates orders of magnitude more support than increasing the dollar amount claimed, a finding with implications for psychology and communication as wel.  We back these findings with an observational design showing that press releases increase constituent support for representatives even after controlling for the amount a legislator actually delivers to the district.  This is done by examining within-respondent differences in support for senators with different rates of credit-claiming, measured by applying ensemble supervised machine learning methods to classify a new collection of over 170,000 House press releases.  We build on these findings in our book, \emph{The Impression of Influence: How Legislator Communication and Government Spending Cultivate a Personal Vote}, which shows how legislators exploit constituents' limited knowledge of appropriations to get credit for bureaucratic spending over which they have little actual control.  

%the complicated federal appropriations process creates ambiguity about who constituents should reward for particularistic spending.  Using original collections of credit claiming messages and new statistical models for text, we show how legislators, with the help of bureaucratic agencies, exploit this ambiguity to create the impression of influence over expenditures. We then document the effect of legislators' strategic credit claiming with a variety of experiments that demonstrate how legislators' deceptive language and constituents' limited knowledge help legislators receive credit for spending and, in turn, cultivate a personal vote. Our findings help explain long standing puzzles about the structure of federal spending, how particularistic spending leads to a personal vote, and Congressional-bureaucratic interactions.  

The next phase will involve an analysis of representative-constituent communication that takes place in social media.  I have gathered a data set comprising of every Twitter communication (``Tweet'') sent to and from every congressional representative since Twitter was founded. I have assembled a similar data set at Facebook, and conducted preliminary analysis released as blogposts.  This is particularly exciting because conventional email and phone communication between representatives and constituents have never been publicly released by representatives, so scholars have very little idea what constituents and representatives discourse about.  Among other things, the project will examine the extent to which representatives claim credit for particularistic benefits versus staking out policy positions versus taunting and/or negatively portraying the opposing party in social media.  The project will also examine how traditional media shape the issues raised in constituent communication, and under what conditions social media activity by representatives shapes media content.  

\section*{Estimating Heterogenous Treatment Effects}
Critical to much of this work is the fact that people respond differently to different messages delivered by different people---heterogenous effects.  While many important discoveries have been made in political science and other disciplines by using randomized experiments to estimate average treatment effects, we often care more about which among an array of treatments has the strongest effect, and how these effects vary across relevant sub-populations.  In ``Estimating Heterogenous Treatment Effects and the Effect of Heterogenous Treatments with Ensemble Methods,'' (working title) we show the close relationship between out-of-sample prediction and accurate estimation of heterogeneous treatment effects, which suggests that cross-validation is a superior method to select from among a large potential set of models and model specifications.  We then show that using out-of-sample, cross-validated performance weight an ensemble of estimates from an array of individual models outperforms other individual methods to estimate heterogenous effects.  We apply this method to two web experiments, illuminating how constituents reward and punish legislators for particularistic spending.  

\section*{Additional Research}
My work with web experiments and large behavioral data sets is part of a broader research agenda focused on applying principles from data science to answer substantive political science questions.  In ``Quantifying Social Media's Political Space,'' (with Robert Bond, forthcoming, \emph{APSR}), we develop a model to simultaneously estimate the ideology of politicians and their supporters, which is strongly correlated with both DW-NOMINATE and survey measures of ideology.  We uses this measure to examine political polarization, how ideology is correlated with our social relationships in society, and how the structure of ideological ties is (negatively) associated with participation rates.  We plan to expand this project to estimate ideology in a cross-national context.  

In ``Bias in the Flesh: Skin Complexion and Stereotype Consistency in Political Campaigns'' (with Ethan Plaut and Maria Jabon, \emph{forthcoming at Public Opinion Quarterly}) we extend findings linking skin complexion to negative stereotypes and adverse real-world outcomes to the realm of political ad campaigns, in which skin complexion can be easily manipulated in ways that are difficult to detect.  We devise a method to measure how dark a candidate appears in an image, and use it to examine how complexion varies with ad content during the 2008 campaign, finding that darker images were more frequent in negative ads---especially in those linking Obama to crime---which aired more frequently as election day approached.  We then conduct an experiment to document how these darker images can activate stereotypes, which shows that a subtle darkness manipulation is sufficient to activate the most negative stereotypes about Blacks---even when the candidate is a famous counter-stereotypical exemplar---Barack Obama.  We then collect further evidence of an evaluative penalty for darker skin based on an observational study measuring affective responses to depictions of Obama with varying skin complexion, presented via the Affect Misattribution Procedure in the 2008 American National Election Survey.  The work demonstrates that darker images are used in a way that compliments ad content, and shows that doing so can negatively affect how individuals evaluate candidates and think about politics.

The next step in this research agenda is to examine the relationship between skin complexion and electoral outcomes.  Toward that end, I've assembled an images database of each primary and general election candidate for national and some state office in 2010 and am merging this with other candidate data.  Preliminary evidence suggests a skin complexion penalty at the ballot box---the measure of skin complexion referenced above is strongly correlated with electoral victory among 2010 minority candidates for national and state-level office.  Future work will expand this database to other election years and examine implications for representation and policy making.  %Furthermore, in experiments conducted with Shanto Iyengar, we show that presenting participants with darker images of a minority candidate decreases feelings of warmth and the probability of voting for the candidate in online experiments, especially for individuals who test high for negative implicit associations with Blacks.

%\section*{Fit with NYU}
%My research agenda compliments existing political science research at NYU, particularly the Social Media and Political Participation project, whose agenda includes formulating and testing theory regarding the cognitive and motivational effects of social media and delivering empirical evidence about the causal mechanisms through which social media affects political participation.  At the same time, the methodological component of my research agenda will compliment the broader political science and data science community, which includes applying machine learning to estimate causal effects with greater precision than other existing techniques, using network structure to generate new measures of ideology, extracting useful features from unstructured text and images using computational and statistical approaches, and employing online labor markets like Mechanical Turk to create new sources of previously unavailable or prohibitively expensive data.


\end{document}









